\documentclass[11pt]{ctexart}
\usepackage{amsmath,amssymb,longtable}
\begin{document}
\section*{P158: 当前完整数学模型（2026-10-20）}
\section*{修订依据与边界}完整模型以OD守恒补齐聚合递推，登乘量允许所有满足类队列/容量的分配；比例分配与有限邻域仅为找可行点的算法。独立检查现有11个启发式解，不要求全局最优；不得把启发式值写作最优值。

\section*{共同计划与时刻表}
方向$x\in\{0,1\}$，当地站序$i=0,\ldots,n-1$，服务$r$，情景$\psi$及概率$E_\psi$。
短交路范围$[\alpha^x,\beta^x]$、行驶$c_i^x$、停站$\phi_i^x$、车底$R$、容量$C$、
到达率$\lambda_{ij}^x(t,\psi)$、权重及归一化均为完整实例数据。
决策$d^x_{r0}$、短交路二元$\eta^x_r$、折返二元$\kappa$及各情景全部乘客流。
\[
L_r^x\le d_{r0}^x\le U_r^x,\quad
d_{ri}^x=d_{r0}^x+\sum_{j<i}c_j^x+\sum_{1\le j\le i}\phi_j^x,\quad
a_{ri}^x=d_{ri}^x-\phi_i^x,\quad
h_{\min}\le d_{ri}^x-d_{r-1,i}^x\le h_{\max}.
\]
每服务最多一个前驱及后继；$\kappa=1$仅可连接不同方向、相同服务类型及同一物理终点的服务，
其折返间隔在$[z_{\min},z_{\max}]$。两方向物理站号映射为$i\leftrightarrow n-1-i$。
\[
RSN=|R^0|+|R^1|-\sum\kappa\le R.
\]
若短交路首端无前驱且非车场则空驶从0开始；若末端无后继且非车场则空驶至$n-1$。
空驶不载客，时间仍按完整区段偏移。折返必须在该物理结束/开始点满足间隔。
\section*{完整OD守恒与选择}
每个OD有目的地索引的三类非负队列：直接、等待长交路、愿意短交路换乘；
初始全零。服务$r$在站$i$之前的新到达为
\[
G^x_{rij}(\psi)=\int_{d^x_{r-1,i}}^{d^x_{ri}}\lambda^x_{ij}(t,\psi)\,dt,\qquad
d^x_{-1,i}=0.
\]
分段常数到达率在数据时域外为零。$\alpha\le i<j\le\beta$为直接类，
其余为II类。短交路$r$在短区间内面对下一个长交路$r'$时，
\[
\xi^x_{ri}={\varkappa_2\over1+\exp[\varkappa_1(\theta^0_{ri}-\theta^0_{r'i})]}.
\]
$\theta^0_{ri}=V^0_{ri}/C$为本站完成下客、尚未上客时的载客率；
后续全长列车的$\theta^0_{r'i}$也采用同一定义。
各时段输入OD人数除以该时段长度得到每分钟到达速率$\lambda$。
若没有后续长交路或本站/本列不满足该条件，$\xi=0$。
II类新需求按$1-\xi,\xi$加入长交路专等/换乘类。外部直接需求加入直接队列。
目的地到达的乘客下车离开系统；短交路在$\beta$强制将目的地更远的乘客卸入
本站保留原目的地的长交路队列，且仅入队一次。所有其他在车乘客继续保留原目的地。
当前服务的合格类别：长交路三类均可，短交路在$[\alpha,\beta)$仅直接和换乘类，其他站不载客。
令$Q^c_{rij}$为加入外部到达和换乘后的类目的地候车量，$V^0_{ri}$为下客后车上量。
\[
B_{ri}=\min\{C-V^0_{ri},\sum_{c,j\text{合格}}Q^c_{rij}\},\quad
0\le B^c_{rij}\le Q^c_{rij},\quad
\sum_{c,j\text{合格}}B^c_{rij}=B_{ri}.
\]
不合格类别的$B^c=0$，剩余队列$Q^c-B^c$保留至下列车；
上车按原目的地增加车上流，故$V_{ri}=V^0_{ri}+B_{ri}\le C$。
未登乘和末班后的到达全保留在期末队列，不删除、不伪称已完成。
这些约束允许全部可行类别分配，\emph{不要求比例分配}。
\section*{完整目标与求解性质}
\[
Z_1(\psi)=\sum_{x,r,i<n-1}[\max(0,V_{ri}/C-\bar\Phi)
+\delta'\max(0,\underline\Phi-V_{ri}/C)],
\quad Z_2(\psi)=\sum_{x,r,i<n-1}(\sum_{\text{合格}c,j}Q^c_{rij}-B_{ri}),\quad Z_3=RSN.
\]
\[
\min\sum_\psi E_\psi\,
{\delta_1Z_1/J_1+\delta_2Z_2/J_2+\delta_3Z_3/J_3\over\delta_1+\delta_2+\delta_3}.
\]
所有情景共用时刻、交路及车底循环。完整模型不固定末班长交路、不限制邻域或迭代次数；
现有算法通过比例登乘及有限邻域构造该模型的可行点。其11个结果仅为校验通过的启发式可行值，
结果为通过ODE物理守恒检查的可行方案。

\end{document}
